\section{Conclusion}

Phase 2 successfully extends the Campus Space Management System for the operating conditions observed during the pilot. The schema now distinguishes advisory maintenance from out-of-service work, preserves maintenance-related booking notifications, normalizes roles and space types, unifies automatic and staff decisions, and migrates Phase 1 records without discarding their history. Five operational triggers automate new-request processing and maintenance notifications while protecting approvals, decided schedules, and usage-session eligibility.

The concurrency experiment demonstrates the central correctness risk and its prevention. The unsafe schedule produces two overlapping approvals, whereas the protected \texttt{BOOKING\_REQUEST}--then--\texttt{SPACES} update-lock protocol leaves exactly one approved request and no overlapping approved pair. The approval-validation trigger extends this protection to direct decision writes that do not originate from the normal stored-procedure path.

The deterministic generator supplies three academic years and 100,000 new booking records. The four required reports cover utilization hours, weekday/hour distribution, suitable-space discovery, and maintenance-escalation impact. Six measured indexes improve every tuned workload; most notably, the room finder reduces logical reads by 97.29\% and elapsed time by 98.04\%.

Finally, the functional-dependency analysis shows that all eleven relations satisfy at least Third Normal Form. The resulting design provides a consistent basis for concurrent operations, operational reporting, and future growth while retaining clear trade-offs between read performance and index-maintenance cost.
